\documentclass[french]{book}
\usepackage{teach}
\usepackage{verbatim}
\usepackage{amsmath,amsfonts,amssymb}
\usepackage{pgfplots}
\usepackage{tikz}
\usepackage{mwe}
\RequirePackage{graphicx}
\RequirePackage{ifpdf}
 \ifpdf
   \DeclareGraphicsRule{*}{mps}{*}{}
 \fi




\RequirePackage{fancyhdr}
\RequirePackage{titlesec}

  \usepackage{fancyhdr}

  \usepackage{amsmath}
  \usepackage{amssymb}
  \usepackage{amsfonts}
  \usepackage{amsthm}
  \usepackage{mathrsfs}

  \usepackage{array}
  \usepackage{multicol}
  \setlength{\multicolsep}{3pt}

  \usepackage{graphicx}
  \graphicspath{{Figures/}}
  \usepackage{pstricks,pst-plot,pst-text,pst-tree,pst-eucl}
  \usepackage[all]{xy}

  \usepackage{fancybox}
  \usepackage{color}
    \usepackage{stmaryrd}
\usepackage{etex}
%\usepackage{preambuleTrm}
%\usepackage{enumitem}
%\usepackage{tikz}
%
\usetikzlibrary{calc}
%\usepackage[xcas]{pro-courbes}
%\usepackage[autolanguage]{numprint}
\usepackage[french,vlined]{algorithm2e}
\usepackage{cclicenses}
%\usepackage{cclicence}
%\usepackage{docacompleter}

\newcommand{\cacheb}[1]{#1}
\newcommand{\cacheg}[1]{#1} 
\newcommand{\cacher}[1]{#1}
\newcommand{\cachegg}[1]{#1} 

% \newcommand{\cacheb}[1]{\dotfill\textcolor{white}{#1}\textcolor{white}{#1}}
% \newcommand{\cacheg}[1]{\dotfill\textcolor{0.9white}{#1}\textcolor{0.9white}{#1}} 

% \newcommand{\cachegg}[1]{\dots\dots\textcolor{0.8white}{#1}\textcolor{0.8white}{#1}} 
\newcolumntype{A}{>{\arraybackslash}X}

\usepackage{textcomp}
\usepackage{mathtools}
%\usepackage[xcas]{pro-statdiv}
%  \usepackage[frenchb]{babel}
  \usepackage{hyperref}

\usepackage{subfig}
\pgfplotsset{compat=newest}
\usepackage[all]{xy}
%\usepackage{geometry}
%\geometry{top=1cm, bottom=1cm, left=2cm, right=2cm}
\usepackage[utf8]{inputenc}
\usepackage[french]{babel}
\redefinecolor{thm}{title}{bgcolor}{alizarin}
\redefinecolor{prop}{title}{bgcolor}{meatbrown}
\redefinecolor{defn}{title}{bgcolor}{olivine}
\definecolor{alizarin}{rgb}{0.82, 0.1, 0.26}
\definecolor{meatbrown}{rgb}{0.9, 0.72, 0.23}
\definecolor{olivine}{rgb}{0.6, 0.73, 0.45}
\usepackage{mathrsfs}
%%
\newcommand{\R}{\mathbb {R}}
%\newcommand{\C}{\mathds {C}}
\newcommand{\N}{\mathbb {N}}
\newcommand{\Z}{\mathbb {Z}}
\newcommand{\Q}{\mathbb {Q}}
\newcommand{\D}{\mathbb {D}}
%%
\newcommand{\se}{\geq}
\newcommand{\ie}{\leq}
\newcommand{\tm}{\times}
\newcommand\binomialCoefficient[2]{%
    % Store values 
    \c@pgf@counta=#1% n
    \c@pgf@countb=#2% k
    %
    % Take advantage of symmetry if k > n - k
    \c@pgf@countc=\c@pgf@counta%
    \advance\c@pgf@countc by-\c@pgf@countb%
    \ifnum\c@pgf@countb>\c@pgf@countc%
        \c@pgf@countb=\c@pgf@countc%
    \fi%
    %
    % Recursively compute the coefficients
    \c@pgf@countc=1% will hold the result
    \c@pgf@countd=0% counter
    \pgfmathloop% c -> c*(n-i)/(i+1) for i=0,...,k-1
        \ifnum\c@pgf@countd<\c@pgf@countb%
        \multiply\c@pgf@countc by\c@pgf@counta%
        \advance\c@pgf@counta by-1%
        \advance\c@pgf@countd by1%
        \divide\c@pgf@countc by\c@pgf@countd%
    \repeatpgfmathloop%
    \the\c@pgf@countc%
}
\makeatother
%\newcommand{\ell}{l}
%\newcommand{\ell'}{l'}
%%
\newcommand{\limn}{\lim_{n\rightarrow +\infty}}
\newcommand{\limo}[1][x]{\ds\lim_{#1\rightarrow 0}}
\newcommand{\limite}[2][x]{\ds\lim_{#1\rightarrow #2}}
\newcommand{\limpinf}[1][x]{\ds\lim_{#1\rightarrow +\infty}}
\newcommand{\limminf}[1][x]{\ds\lim_{#1\rightarrow -\infty}}
\newcommand{\limiteg}[2][x]{\ds\lim_{#1\xrightarrow{<}#2}}
\newcommand{\limited}[2][x]{\ds\lim_{#1\xrightarrow{>}#2}}
\newcommand{\anglevec}[2]{(\overrightarrow{#1};\overrightarrow{#2})}
\newcommand{\co}[2]{C_#2^#1}
\newcommand{\ssi}{\Longleftrightarrow}
%%
\newcommand{\vect}[1]{\overrightarrow{#1}}
\newcommand{\oij}{(\textit{O},{\overrightarrow{i}},{\overrightarrow{j}})}
%
\newcommand{\cp}[2]{%
        \begin{pmatrix}
        #1\\
        #2
        \end{pmatrix}%
        }
%
\newcommand{\calb}{\mathscr{B}}
\newcommand{\calc}{\mathscr{C}}
\newcommand{\cald}{\mathscr{D}}
\newcommand{\cale}{\mathscr{E}}
\newcommand{\calf}{\mathscr{F}}
\newcommand{\calp}{\mathscr{P}}
\newcommand{\fr}[2]{\dfrac{#1}{#2}}
\newcommand{\cals}{\mathscr{S}}
\newcommand{\calt}{\mathscr{T}}
\newcommand{\ve}[1]{\overrightarrow{#1}}
\newcommand{\pa}[1]{(#1)}
\newcommand{\ps}[2]{\overrightarrow{#1}.\overrightarrow{#2}}
\newcommand{\bbr}{\mathbb{R}}
\newcommand{\al}{\alpha}
\newcommand{\la}{\lambda}
\newcommand{\DR}{\cald}
\newcommand{\PR}{\calp}
\newcommand{\nor}[1]{ \Vert #1 \Vert}
\newcommand{\abv}[1]{\vert #1 \vert}
\newcommand{\coordpp}[3]{\begin{pmatrix}
#1\\
#2\\
#3
\end{pmatrix}}
%
\newcommand{\suite}[1][u]{\left( #1_{n} \right)_{n\in\N}}
\newcommand{\suitep}[1][u]{\left( #1_{n} \right)_{n\in\N^{*}}}
\usetikzlibrary{arrows}


\begin{document}
\tableofcontents
%Nous  allons développer  dans cette  leçon une  des grandes
%nouveautés  de terminale  : le  lien  entre algèbre  et géométrie,  deux
%domaines jusqu'ici  étrangers. L'outil central sera  le produit scalaire
%qui  sera  donc étudié  avec  attention.  Vous  verrez (peut-être)  l'an
%prochain qu'un \og 
%ensemble  de  vecteurs\fg{}  dans  lequel  on peut  définir  un  produit
%scalaire est appelé espace euclidien. 

        \chapter{Géométrie euclidienne} 




\section{Construisons le produit scalaire}\label{ps}\index{produit scalaire |( }

\subsection{À la recherche d'une définition}

Vous avez remarqué au moment de l'étude de la fonction exponentielle que
son  mode de  construction  n'était  pas unique  :  certains partent  de
l'équation différentielle  $y'=y$, d'autres de  l'équation fonctionnelle
$f(x+y)=f(x)\times f(y)$, etc.
\begin{comment}  d'autres encore de la suite  de terme général
$(1+\fr{x}{n})^n$               ou               même              de
$\sum_{n=0}^{+\infty}\fr{x^n}{n!}$...
\end{comment}

Pourtant, tous construisent la même
fonction,  car  ces définitions  s'avèrent  équivalentes.  La nature  du
problème  étudié,  le  niveau  d'enseignement,  etc. incitent  alors  à
préférer l'une ou l'autre des définitions. \\

C'est encore le cas pour le produit scalaire. Nous choisirons la méthode
la  plus adaptée  au  programme  de terminale  et  n'utiliserons que  le
théorème de  Pythagore en  admettant que l'on  peut munir  l'Espace d'un
repère orthonormé. 

\begin{defn}[ Produit scalaire dans l'Espace]
 Soient $\ve{u}$ et $\ve{v}$ de coordonnées
respectives $(x,y,z)$ et $(x', y', z')$  dans un repère orthonormé
de l'Espace.\\

On appelle produit scalaire de $\ve{u}$ et $\ve{v}$  que l'on note
$\ve{u}\cdot\ve{v}$ le réel

$$\ve{u}\cdot\ve{v} = xx'+yy'+zz'$$

\end{defn}


Cette définition est déroutante : elle dépend du choix d'un repère : y aurait-il un produit scalaire par repère, ce qui pourrait s'avérer fort gênant~! Nous allons en fait nous assurer que cette définition en est bien une, c'est à dire qu'elle ne dépend pas du choix du repère.\\

Tout d'abord, un petit dessin :

\begin{center}
\includegraphics{../Figures/geo3d.1} 
\end{center}

Soit $\ve{OM}$ un représentant du vecteur $\ve{u}$. Nous pouvons
calculer  $OM^2$, c'est à dire le carré de  la norme du vecteur
$\ve{u}$ en fonction de $x$, $y$ et $z$.\\

D'après le  théorème de Pythagore :\\

$OM^2=OM'^2+z^2$ et d'autre part, $OM'^2=x^2+y^2$. Finalement on
obtient que $OM^2=x^2+y^2+z^2$.

Cette formule permet surtout de
prouver, avec un minimum de sens de l'observation, que



\begin{prop}[ Norme et produit
scalaire]
$$\ve{u}\cdot\ve{u}=\nor{\ve{u}}^2$$
\end{prop}



Vous aurez remarqué qu'il n'est pas fait mention de repère dans
cette propriété et qu'elle est donc indépendante du repère choisi.

Intéressons-nous maintenant à $\nor{\ve{u}+\ve{v}}^2$

\begin{eqnarray*}
\nor{\ve{u}+\ve{v}}^2 &=& (x+x')^2+(y+y')^2+(z+z')^2\\
                    &=&
                    x^2+y^2+z^2+x'^2+y'^2+z'^2+2(xx'+yy'+zz')\\
                    &=&
                    \nor{\ve{u}}^2+\nor{\ve{v}}^2+2(\ve{u}\cdot\ve{v})
\end{eqnarray*}

Je vous laisse le soin d'obtenir une formule similaire en
remplaçant $\ve{v}$ par son opposé.



\begin{prop}[Expression du produit scalaire en fonction de la norme]\label{psnorme}
  \begin{eqnarray*} \ve{u}\cdot\ve{v} &=&
\fr{1}{2}\pa{\nor{\ve{u}+\ve{v}}^2-\nor{\ve{u}}^2-\nor{\ve{v}}^2}\\
&=& \fr{1}{2}\pa{\nor{\ve{u}}^2+\nor{\ve{v}}^2-\nor{\ve{u}-\ve{v}}^2}
\end{eqnarray*}
\end{prop}




Ces expressions peuvent paraître déroutantes car très compliquées, voire
inexploitables. \\
 Cependant, elles  offrent  \textbf{ un énorme  avantage} :  elles
prouvent que le  produit scalaire ne dépend pas  des coordonnées et donc
de  la  base  choisie.  Nous  pouvons  donc  utiliser  sans  retenue  la
définition qui  en est  bien une finalement.

Notez de  plus que cette dernière  formulation est valable  dans le Plan
comme  dans  l'Espace  car  la  démonstration  est  indépendante  de  la
\og dimension\fg{}    choisie,  ce   qui  améliore   encore   son  \og champs
d'action\fg{} .  



Toujours  grâce  à  cette   propriété,  nous  pouvons  prouver  que  les
propriétés suivantes sont vraies quelque soit la base choisie 

\begin{prop}[Propriétés de bilinéarité et de
symétrie]
\begin{itemize}
\item $\ve{u}\cdot\ve{v}=\ve{v}\cdot\ve{u}$ 
\item $(\pa{\alpha\ve{u}}\cdot\ve{v}=\ve{u}\cdot\pa{\al\ve{v}}=\al\pa{\ve{u}\cdot\ve{v}}\qquad\text{avec
}\quad \al\in\bbr$
\item $\ve{u}\cdot\pa{\ve{v}+\ve{w}}=\ve{u}\cdot\ve{v}+\ve{u}\cdot\ve{w}$
\end{itemize}
\end{prop}


\newpage

\subsection{Produit scalaire et cosinus}

Nous   allons   maintenant   interpréter  géométriquement   le   produit
scalaire.   Considérons   deux   vecteurs   $\ve{u}$  et   $\ve{v}$   de
représentants $\ve{OA}$ et $\ve{OB}$ 



\begin{center}
\includegraphics{../Figures/geo3d.2} 
\end{center}


Nous  allons   essayer  de  prouver   le  résultat  vu  en   première  :
$\ps{u}{v}=OH\times OB$ à partir de la propriété sur les normes. \\

Le thèorème de Pythagore permet d'écrire d'une part que $OH^2+AH^2=OA^2$ et d'autre part que $(OB-OH)^2+AH^2=AB^2$. \\ 

On élimine $AH^2$ pour obtenir 
$$                                                               OB\times
OH=\fr{1}{2}\pa{OA^2+OB^2-AB^2}=\fr{1}{2}\pa{\nor{\ve{u}}^2+\nor{\ve{v}}^2-\nor{\ve{u}-\ve{v}}^2}=\ve{u}\cdot\ve{v} $$ 

\textit{(On obtiendrait d'ailleurs un résultat similaire si $H$ était \og de
l'autre côté\fg{}  de $O$, en transformant $OH$ en $-OH$, puisque
dans ce cas $HB=OB+OH$.)}\\




Ici un problème se pose : on sait que $OB=\nor{\ve{v}}$ mais comment
interpréter $OH$  ? On a  envie d'écrire que  $OH=\nor{\ve{u}}\times \cos
\theta$ comme  au collège, mais  qu'est-ce que $\theta$ ? \emph{(Pour rappel nous somme dans l'Espace et non plus dans le Plan)}. Qu'est-ce que $\cos\theta$~?  \\


Nous allons trancher dans le vif. Considérons deux vecteurs non nuls $\ve{u}$ et $\ve{v}$. Notons 
$$ \cos\anglevec{u}{v}=\fr{\ps{u}{v}}{\nor{\ve{u}}\times\nor{\ve{v}}} $$

La définition du produit scalaire à partir des normes étant la
même dans le Plan et l'Espace, on va pouvoir \textbf{étendre} la notion d'angle et parler de l'angle entre deux vecteurs de l'espace à partir de la même formule.\\

On peut noter que deux vecteurs définissant un plan \emph{(vectoriel)},
il est naturel d'étendre la notion d'angle de vecteurs du Plan à
l'Espace.\\

Demeure un dernier problème : le Plan était \underline{orienté}, or nous
n'avons pas parlé d'orientation de l'Espace (et nous ne le ferons
pas cette année). Les angles de vecteurs de l'Espace seront donc
définis \underline{au signe près} et mesurés dans $[0,\pi]$, ce qui rejoint la
notion d'angle géométrique vu au Collège.



\begin{thm}[Produit scalaire et cosinus]\label{pscos}
  Soient $\ve{u}$ et $\ve{v}$ deux vecteurs non nuls de l'Espace. On a
$$ \cos\anglevec{u}{v}=\fr{\ps{u}{v}}{\nor{\ve{u}}\times\nor{\ve{v}}} $$
avec $ \anglevec{u}{v}$ l'unique antécédent de
$\fr{\ps{u}{v}}{\nor{\ve{u}}\times\nor{\ve{v}}}  $ par la
fonction cosinus dans $[0,\pi]$.
\end{thm}

On retrouve de cette manière la formule bien connue $\ps{u}{v}=\nor{\ve{u}}\times\nor{\ve{v}}\times\cos\theta$.


\subsection{Produit scalaire et orthogonalité}\index{orthogonalité}

Nous allons maintenant aborder la notion qui a motivé
l'introduction du produit scalaire en Terminale, à savoir
\underline{l'orthogonalité de deux vecteurs}.
%, et nous allons, par la même
%occasion, être confrontés une nouvelle fois au problème de la
%poule et de l'\oe uf.


 Considérons la figure suivante :

\begin{center}
\includegraphics{../Figures/geo3d.3} 
\end{center}


Si $(AB)\bot(BC)$, alors le théorème de Pythagore assure que
$\nor{\ve{u}+\ve{v}}^2=\nor{\ve{u}}^2+\nor{\ve{v}}^2$.\\

Or $\ve{u}\cdot\ve{v} =
\fr{1}{2}\pa{\nor{\ve{u}+\ve{v}}^2-\nor{\ve{u}}-\nor{\ve{v}}}$, donc
$\ve{u}\cdot\ve{v}=0$.\\



Inversement, si $\ve{u}\cdot\ve{v}=0$, alors
$\nor{\ve{u}+\ve{v}}^2=\nor{\ve{u}}^2+\nor{\ve{v}}^2$ et la réciproque
du théorème de Pythagore assure que $(AB)\bot (AC)$.\\

Finalement on a le théorème suivant:


\begin{thm}[Produit scalaire et orthogonalité]
  Soient deux vecteurs $\ve{u}$ et $\ve{v}$ non nuls et trois points
$O$, $A$ et $B$ tels que $\ve{u}=\ve{OA}$ et $\ve{v}=\ve{OB}$. Les
trois propositions suivantes sont équivalentes :

\begin{itemize}
\item[$\bullet$] $(OA)$ et $(OB)$ sont perpendiculaires
\item[$\bullet$] $\ve{u}\cdot\ve{v}=0$
\item[$\bullet$] $\ve{u}$ et $\ve{v}$ sont orthogonaux. On notera $\ve{u}\bot\ve{v}$.
\end{itemize}

Par souci de cohérence, on dira que le vecteur nul est orthogonal
à tout vecteur.
\end{thm}


Vous noterez que ces résultats sont cohérents avec la définition
du cosinus.

%\index{produit scalaire |) }
\section{Applications du produit scalaire}

\begin{comment}
\subsection{Affine $ \leftrightharpoons $ Vectoriel}

Avant  d'aller plus  loin, commençons  par parler  des limites  de notre
exposé. Dans  l'enseignement secondaire, vous n'avez  travaillé que dans
des  espaces \textsl{affines},  c'est  à dire,  grossièrement, avec  des
ensembles  de  points~:~droites, cercles,  segments  de droites,  angles
géométriques,  etc. Or  vous découvrirez  l'an prochain  la  géométrie à
partir d'espaces \textsl{vectoriels},  c'est à dire, grossièrement, avec
des  ensembles stables  par  combinaison linéaire  comme l'ensemble  des
vecteurs du Plan~:~si $\ve{u}$ et  $\ve{v}$ sont des vecteurs du Plan et
$\la$  et $\mu$  des réels,  alors $\la\ve{u}+\mu\ve{v}$  est  encore un
vecteur du  plan. Ensuite, vous  définirez les espaces affines  à partir
des  espaces vectoriels\footnote{Ces  structures dépassent  largement la
  Géométrie~:~vous découvrirez par  exemple que l'ensemble des solutions
  d'une  équation différentielle  linéaire du  second ordre  sans second
  membre a une structure d'espace vectoriel de dimension 2 et l'ensemble
  des solutions  d'une équation différentielle linéaire  du second ordre
  avec  second membre  a une  structure d'espace  affine de  dimension 2
  comme un plan dans l'Espace \og géométrique\fg{}  du lycée.} 

Nous ne pouvons bien sûr  pas suivre ce cheminement. Nous nous appuirons
donc  sur les  résultats \og affines\fg{}   vus au  lycée et  au  collège et
n'évoquerons les espaces vectoriels qu'à titre d'exercice. 

Il est malgré  tout très important de bien  distinguer ces deux domaines
et  nous  ferons souvent  des  appels  intuitifs  à ces  notions.   Ayez
cependant bien  en tête que  nous travaillerons principalement  dans des
espaces affines  et que le recours  aux outils vectoriels  doit se faire
avec prudence. 

\end{comment}

\subsection{Droites orthogonales}

\begin{defn}[Droites orthogonales]
  Soit $\DR$ une droite de vecteur directeur $\ve{u}$ et $\Delta$ une droite de vecteur directeur $\ve{v}$. \\
On dira que les droites  $\DR$ et $\Delta$ sont \textbf{orthogonales} et
on notera $\DR\bot\Delta$ si  et seulement si $\ve{u}\bot\ve{v}$ c'est à
dire si et seulement si $\ps{u}{v}=0$
\end{defn}

\begin{center}
\includegraphics{../Figures/geo3z.1} 
\end{center}




\subsection{Droites et plans orthogonaux}


\begin{defn}[Orthogonalité d'une droite et d'un plan]
  Soit $\DR$ une droite de vecteur directeur $\ve{u}$ et $\PR$ un plan de vecteurs directeurs $\ve{v_1}$ et $\ve{v_2}$.\\
On dira que $\DR$ et $\PR$ sont orthogonaux et on notera $\DR\bot\PR$ si et seulement si $\ve{u}\bot\ve{v_1}$ et $\ve{u}\bot\ve{v_2}$ c'est à dire si et seulement si $\ps{u}{v_1}=0$ et $\ps{u}{v_2}=0$
\end{defn}

\begin{center}
\includegraphics{../Figures/geo3z.2} 
\end{center}

Notez bien  que, comme dans le  cas des droites,  l'orthogonalité est en
fait  une  propriété  vectorielle.  On  peut  néanmoins  en  donner  une
formulation équivalente en faisant intervenir des points~: 

\begin{prop}[]
  $\DR\bot\PR$ si et  seulement si, pour tous points M  et N de $\PR$
  on a $\ps{u}{MN}=0$
\end{prop}

En  effet, $\ve{v_1}$  et $\ve{v_2}$  étant des  vecteurs  directeurs de
$\PR$,    il    existe   deux    réels    $x$    et    $y$   tels    que
$\ve{MN}~=~x~\ve{v_1}~+~y~\ve{v_2}$ et donc 
$$ \ps{u}{MN}=\ve{u}\cdot\pa{x\ve{v_1}+y\ve{v_2}}=x\ps{u}{v_1}+y\ps{u}{v_2}=0 $$

Notez qu'on  retrouve la  définition 3 en  prenant $x=1$ et  $y=0$, puis
$x=0$ et $y=1$. 

\subsection{Vecteur normal à un plan}\index{vecteur normal} 

La direction d'une droite est  déterminée par la donnée d'un vecteur. Le
problème pour un  plan, c'est qu'on a besoin a  priori de deux vecteurs,
ce qui  est doublement plus compliqué  et qui gène  le mathématicien qui
vise  toujours la  simplicité. 
C'est  ici  que l'intuition vectorielle vient  nous sauver~:~si  un plan a  deux directions,  il n'a
\textbf{qu'une seule  \og direction orthogonale\fg{}} et  donc la direction  du plan
sera \textbf{entièrement determinée}  par la donnée d'un vecteur  orthogonal à la
direction du plan, qu'on appelera \textbf{vecteur normal} au plan.  

%Nous ne nous  lancerons pas dans une preuve  artificielle à notre niveau
%et qui sera  simple avec les outils que  vous découvrirez l'an prochain,
%mais nous nous contenterons de cette approche intuitive. 

\begin{defn}[Vecteur  normal à un plan]
  \'Etant donné  un plan $\PR$, on
  appelera \textbf{vecteur normal} à  $\PR$ tout vecteur directeur d'une
  droite orthogonale à $\PR$
\end{defn}

%je n'ai toujours pas retrouvé mon ordinateur portable






\begin{center}
\includegraphics{../Figures/geo3z.3} 
\end{center}


Nous admettrons alors les résultats suivants :


\begin{thm}[Détermination d'un plan à l'aide d'un vecteur normal]
  Soit $\PR$ un plan, A un point de $\PR$ et $\ve{n}$ un vecteur normal à $\PR$, alors le point M appartient à $\PR$ si et seulement si $\ps{AM}{n}=0$.\\


Réciproquement, si A est un point quelqconque et $\ve{n}$ un vecteur non nul, alors l'ensemble des points M tels que $\ps{AM}{n}=0$ est un plan.
\end{thm}










\subsection{Distance d'un point à un plan}

\begin{defn}[Distance d'un point à  un plan]
  Soit M un  point, $\PR$ un
  plan.  On appelle distance  de M  à $\PR$  la distance  MH, avec  H le
  projeté orthogonal de M sur $\PR$
\end{defn}


%Vous pouvez  commencer, pour vous  échauffer, par montrer à  l'aide d'un
%dessin  et d'un petit  raisonnement utilisant  un théorème  connu depuis
%fort  longtemps,  que  MH est  en  fait  le  minimum des  distances  MN,
%$N\in\PR$. 



Il reste à exprimer cette distance MH. On suppose connu le plan $\PR$ et
donc  un point  quelconque A  de  $\PR$ et  un de  ses vecteurs  normaux
$\ve{n}$. 

%\break

Aidons-nous alors du dessin suivant :

\begin{center}
\includegraphics{../Figures/geo3z.4} 
\end{center}


Cela  doit   maintenant  devenir  un   réflexe  :  qui   dit  projection
orthogonale, dit  produit scalaire, donc la distance  MH devrait pouvoir
s'exprimer en fonction de $\ps{AM}{n}$. 

Or on obtient successivement

\begin{align*}
\ps{AM}{n}&= \pa{\ve{AH}+\ve{HM}}\cdot\ve{n}\\
                   &= \ps{AH}{n}+\ps{HM}{n}\\
                   &=\ps{HM}{n}\\
                   &=\pm d(M,\PR)\times \nor{\ve{n}}
\end{align*}

Le signe étant choisis selon les sens respectifs de $\ve{HM}$ et $\ve{n}$ \emph{(positif s'il est dans le même sens et négatif sinon) }. Le problème, c'est qu'on ne les connait pas a priori. Mais ce n'est plus un problème si on cache les signes sous une valeur absolue, puisque c'est la distance \emph{ donc positive } qui nous intéresse.



Ainsi, $\abv{\ps{AM}{n}}=d(M,\PR)\times \nor{\ve{n}}$ et donc 


\begin{thm}[Calcul de la distance d'un point à un plan]
  Soit $\PR$ un plan passant par le point A et de vecteur normal $\ve{n}$ et soit M un point quelconque. Alors
$$ d(M,\PR)=\fr{\abv{\ps{AM}{n}}}{\nor{\ve{n}}} $$
\end{thm}


\subsection{Plans parallèles}

Voici une utilisation bien pratique de la notion de vecteur normal~:

\begin{prop}[Plans parallèles]
  Deux plans sont parallèles si et seulement si leurs vecteurs normaux sont colinéaires.
\end{prop}


\begin{center}
\includegraphics{../Figures/geo3z.5} 
\end{center}


\subsection{Plans perpendiculaires}

En voici une autre souvent utile en exercice~:


\begin{prop}[Plans perpendiculaires]
  Deux plans sont perpendiculaires si et seulement si leurs vecteurs normaux sont orthogonaux.
\end{prop}





\begin{center}
\includegraphics{../Figures/geo3z.6} 
\end{center}
Remarquez en particulier que $\ve{n_1}$ et un vecteur directeur de $\PR_2$ et $\ve{n_2}$ un vecteur directeur de~$\PR_1$.

%\break
\newpage
\section{Géométrie analytique}

\subsection{Représentation paramétrique d'un plan}

Soit $\PR$ un plan passant  par $A$ de coordonnées $(x_A,y_A,z_A)$ et de
vecteurs       directeurs       $\ve{u}      \coordpp{a}{b}{c}$       et
$\ve{u'}\coordpp{a'}{b'}{c'}$. Alors un point M de coordonnées $(x,y,z)$
appartient  à $\PR$ si,  et seulement  s'il existe  deux réels  $\la$ et
$\mu$ tels que $$ \ve{AM}=\la\ve{u}+\mu\ve{v} $$ 
c'est à dire, en appliquant cette relation  aux coordonnées

\begin{center}
$$\begin{cases}
x =x_A+\la a+\mu a'\\
y =y_A+\la b+\mu b'\\
z =z_A+\la c+\mu c'
\end{cases}$$
\end{center}

C'est ce qu'on appelle une représentation paramétrique du plan $\PR$. Il
est également  utile de retrouver les éléments  caractéristiques du plan
(~un   point  et   deux  vecteurs   directeurs~)  à   partir   de  cette
représentation. 



\subsection{\'Equation cartésienne d'un plan}

On utilise plus  couramment une équation cartésienne de  plan. En effet,
nous avons vu qu'un plan  était entièrement déterminé par la donnée d'un
point et d'un vecteur normal.  

Soit donc $\PR$  un plan passant par $A$  de coordonnées $(x_A,y_A,z_A)$
et de vecteur normal $\ve{n}\coordpp{a}{b}{c}$.  
Alors M de coordonnées $(x,y,z)$  appartient à $\PR$ si, et seulement si
$\ve{AM}\bot\ve{n}$, c'est à dire 

\begin{center}
\begin{align*}
\ps{AM}{n}&=a(x-x_A)+b(y-y_A)+c(z-z_A)\\
                    &=ax+by+cz-(ax_A+by_A+cy_A)\\
                    &=0
\end{align*}
\end{center}

C'est à dire encore, en posant $d=ax_A+by_A+cz_A$


\begin{prop}[Équation cartésienne d'un plan]
  $$M(x,y,z)\in\PR \ssi ax+by+cz=d \text{ avec  } \ve{n}\coordpp{a}{b}{c}
\text{ un vecteur normal à } \PR $$
\end{prop}






\newpage

\subsection{Représentation paramétrique d'une droite}

Quant aux droites, nous n'avons guère le choix~:~un point M appartient à
la  droite $\DR$  passant par  A  de coordonnées  $(x_A,y_A,z_A)$ et  de
vecteur  directeur $\ve{u}$  de coordonnées  $\coordpp{a}{b}{c}$  si, et
seulement s'il existe un réel $\la$ tel que 

$$ \ve{AM}=\la\ve{u} $$

c'est à dire si, et seulement si

\begin{center}
$ \begin{cases}
x =x_A+\la a\\
y =y_A+\la b\\
z =z_A+\la c
\end{cases}$
\end{center}

Notez  bien que  $\la$  est  le paramètre  de  la représentation  (~d'où
l'appellation~), mais aurait pu porter tout autre nom. 







\subsection{Sphères}


\begin{defn}[Sphère]
  
Une sphère de centre $C$ et de  rayon $r$ est l'ensemble des points M de
l'espace vérifiant~:

$$\Vert \ve{CM} \Vert =r$$
\end{defn}



Analytiquement, dans un repère orthonormé, nous écrirons~:


$$(x-x_C)^2+(y-y_C)^2+(z-z_C)^2=r^2$$





\section{Barycentres}

\subsection{Quelques rappels}




\begin{defn}[Barycentre]
  Soient $A_1,\ A_2\,...,\,A_n$ $n$ points distincts du plan ou de l'espace. Soient $\al_1,\ \al_2\,...,\,\al_n$ $n$ réels tels que $\al_1+\al_2+\cdots+\al_n\neq 0$, alors on appelle barycentre des points $A_1,\ A_2\,...,\,A_n$ affectés des coefficients $\al_1,\ \al_2\,...,\,\al_n$ le point G défini par 
$$\al_1\ve{GA_1}+\al_2\ve{GA_2}+\cdots+\al_n\ve{GA_n}=\ve{0}$$
\end{defn}


\c Ca, vous connaissez par c\oe ur. Il est maintenant utile de connaître la manipulation suivante:\\
 pour exprimer G en fonction de points connus. Soit donc M un tel point (en général O, l'origine du repère), alors

$$\al_1\pa{\ve{GM}+\ve{MA_1}}+\al_2\pa{\ve{GM}+\ve{MA_2}}+\cdots+\al_n\pa{\ve{GM}+\ve{MA_n}}=\ve{0}$$
 puis, finalement, en regroupant astucieusement les termes, on obtient~:


 \begin{thm}[Fonction vectorielle de Leibniz]
   $$\ve{MG}=\fr{1}{\al_1+\al_2+\cdots+\al_n}\pa{\al_1\ve{GA_1}+\al_2\ve{GA_2}+\cdots+
\al_n\ve{GA_n}}$$
 \end{thm}


En particulier, en prenant $M=O$, on obtient 
$$\begin{cases}x_G=\fr{1}{\al_1+\al_2+\cdots+\al_n}\pa{\al_1 x_{A_1}+\al_2x_{A_2}+\cdots+\al_nx_{A_n}}\\
y_G=\fr{1}{\al_1+\al_2+\cdots+\al_n}\pa{\al_1 y_{A_1}+\al_2y_{A_2}+\cdots+\al_n y_{A_n}}\\
z_G=\fr{1}{\al_1+\al_2+\cdots+\al_n}\pa{\al_1 z_{A_1}+\al_2z_{A_2}+\cdots+\al_nz_{A_n}}
\end{cases}$$


\subsection{Quelques nouveautés}

%Nous commenterons en classe les résultats suivants~:


\begin{thm}[Caractérisations barycentriques des droites et des plans]
  \begin{itemize}
  \item[$\bullet$]  La droite (AB) est l'ensemble des barycentres des points A et B.
\item[$\bullet$] Le segment $[A;B]$ est l'ensemble des barycentres de A et B affectés de coefficients positifs.
\item[$\bullet$] Soient A, B et C trois points non alignés. Le plan (ABC) est l'ensemble des barycentres des points A, B et C
  \end{itemize}
\end{thm}



Ce chapitre  est assez riche en  cours, mais surtout  en applications de
ces résultats~:~les  exercices qui suivent peuvent  donc être considérés
comme faisant partie intégrante du cours. 


\end{document}
